\chapter{TCP streams and the socket interface}
\chapterintro{TCP turns unreliable packet delivery into an ordered stream between two endpoints. Sockets are the operating-system interface through which our Python process joins that stream.}

\begin{objectives}
\begin{itemize}
  \item Explain connection establishment, ordered delivery, acknowledgments, and closing.
  \item Distinguish reliability, flow control, congestion control, and application timeouts.
  \item Describe listening and connected sockets and the bind--listen--accept lifecycle.
  \item Predict partial reads, partial writes, blocking, and slow-client behavior.
\end{itemize}
\end{objectives}

\section{A reliable byte stream}

TCP provides two applications with an ordered, bidirectional sequence of bytes.
It numbers bytes conceptually, detects loss, retransmits data, removes duplicates,
and delivers data to the receiving application in order. If a later segment
arrives before an earlier one, the implementation normally waits rather than
showing the application a gap.

Reliability does not mean success is guaranteed. A cable can remain disconnected,
a process can crash, or a peer can stop reading forever. TCP either continues
trying within its policies or reports failure. Applications still need deadlines,
cancellation, and error handling.

\section{Establishing a connection}

The familiar three-way handshake synchronizes connection state and initial
sequence numbers:

\begin{enumerate}
  \item The client sends a segment with the SYN flag.
  \item The server replies with SYN and ACK.
  \item The client acknowledges the server's SYN.
\end{enumerate}

The handshake also negotiates capabilities. It costs a network round trip before
ordinary application data can flow, which is one reason connection reuse matters.
A server's listening socket and backlog participate before application code
accepts the established connection.

\begin{figure}[htbp]
\centering
\begin{tikzpicture}[
  actor/.style={font=\sffamily\bfseries},
  msg/.style={-{Latex[length=2mm]},Accent}
]
\node[actor] (client) {client};
\node[actor,right=105mm of client] (server) {server};
\draw[Rule] (client.south) -- ++(0,-50mm);
\draw[Rule] (server.south) -- ++(0,-50mm);
\draw[msg] ([yshift=-12mm]client.south) -- node[above,font=\small]{SYN} ([yshift=-12mm]server.south);
\draw[msg] ([yshift=-25mm]server.south) -- node[above,font=\small]{SYN + ACK} ([yshift=-25mm]client.south);
\draw[msg] ([yshift=-38mm]client.south) -- node[above,font=\small]{ACK} ([yshift=-38mm]server.south);
\end{tikzpicture}
\caption{A simplified TCP three-way handshake.}
\label{fig:tcp-handshake}
\end{figure}

\section{Sequence, acknowledgment, and retransmission}

A TCP acknowledgment says, roughly, ``I have received all bytes before this
next expected sequence number.'' Senders keep unacknowledged data so it can be
retransmitted. They infer loss through timers and duplicate acknowledgment
patterns. Checksums detect corruption within their scope.

The application does not receive these segments directly. It writes bytes to a
socket buffer and reads bytes from another buffer. TCP can split one write into
many segments or combine data from several writes. Network offload mechanisms
make the relationship even less visible. The application must treat the result
as a stream.

\section{Two kinds of control}

\term{Flow control} protects the receiver. The receiver advertises how much
buffer capacity is available; a sender that ignores this would overwhelm the
destination. \term{Congestion control} protects the network path. The sender
adjusts how much data is in flight based on observed conditions such as loss
and delay.

These mechanisms can make a write block or make progress slowly. A successful
write means bytes were accepted by the local socket machinery, not that a remote
application has already processed them. Likewise, a successful TCP handshake
does not mean an HTTP request will be accepted.

\section{Closing is a protocol too}

Each direction of a TCP connection closes independently with FIN and
acknowledgment exchange. A peer can finish sending while still receiving. An
abrupt failure can produce a reset (RST), which discards the graceful-close
story. Servers must expect clients to disappear between any two operations.

The endpoint that actively closes can remain in a \code{TIME\_WAIT} state. This
helps prevent delayed segments from an old connection being mistaken for a new
one with the same endpoint tuple. It is normal protocol housekeeping, not
automatically evidence of a leak.

\section{Sockets: the application boundary}

A socket is an operating-system object representing a communication endpoint.
The typical server lifecycle is:

\begin{enumerate}
  \item \term{create} a TCP socket;
  \item \term{bind} it to a local address and port;
  \item mark it as \term{listening} with a backlog;
  \item \term{accept} a connection, producing a new connected socket;
  \item receive and send bytes on the connected socket;
  \item close the connected socket and eventually the listener.
\end{enumerate}

The listener does not carry HTTP request bytes for an accepted client. The
connected socket does. This separation lets one accept loop hand connections to
workers while continuing to receive new arrivals.

\subsection{Blocking and timeouts}

A blocking call waits until it can make progress or report an event. Accept can
wait for a client. Receive can wait for bytes. Send can wait for buffer capacity.
Simple blocking code is ideal for learning because control flow is visible, but
unbounded waiting gives a slow or silent peer control over a worker.

A timeout changes an infinite wait into an explicit failure path. The server we
build will set a read/idle timeout, limit header and body sizes, bound the number
of workers and queued connections, and limit requests per persistent connection.
These constraints are part of correctness, not optional tuning.

\begin{securityboundary}
An attacker does not need high bandwidth to exhaust a naive server. Many
connections that send headers one byte at a time can occupy all blocking workers.
Deadlines, byte limits, bounded queues, and graceful rejection reduce this risk.
\end{securityboundary}

\section{Partial operations are ordinary}

A receive call returns up to the requested number of bytes. It may return fewer
because that is what is currently available. An empty byte string indicates
orderly end-of-stream. A low-level send can accept fewer bytes than supplied;
Python's \code{sendall} repeatedly sends until all bytes are accepted or an error
occurs.

The correct question is never ``Did one read return my request?'' It is ``Does
my accumulated buffer contain enough bytes, according to HTTP framing, to parse
one complete request?''

\section{Why not UDP?}

UDP sends independent datagrams and does not provide TCP's ordered reliable
stream. Applications that choose it must accept loss/reordering or implement
the needed behavior. HTTP/3 uses QUIC over UDP, but QUIC itself supplies secure,
reliable streams and congestion control. ``Uses UDP'' does not mean ``ignores
reliability.'' We implement HTTP/1.1 over TCP because its text framing exposes
the fundamentals clearly.

\summarytitle

TCP establishes state, presents ordered bytes, recovers from common packet loss,
and controls sender behavior, but applications still need framing, timeouts,
limits, and failure handling. A listening socket accepts connected sockets; each
connected socket carries an unstructured byte stream from HTTP's perspective.

\reviewtitle
\begin{enumerate}
  \item Why can one application write become several TCP segments?
  \item What is the difference between flow control and congestion control?
  \item What does an empty result from a receive operation mean?
  \item Why does the server need both a listener and connected sockets?
\end{enumerate}

\begin{exercise}[Stream outcomes]
A client writes a 600-byte header and a 200-byte body. List at least four legal
ways a server's sequence of receive calls could divide those 800 bytes. State
which HTTP rule tells the server when the body is complete.
\end{exercise}

\begin{exercise}[Slow client budget]
Suppose a server has 16 blocking workers and no read timeout. Explain how 16
clients can prevent a seventeenth valid client from being handled. Propose a
timeout, queue, and rejection policy and describe the trade-off of each.
\end{exercise}

